%\section{Introduction}

%- Recent structure-property studies have linked spinodal decomposition to stiffening of bcc delta ferrite, this paper sets out a theoretical foundation for the predictability of this phenomenon from quantum elastodynamic first principle, by combining elastodynamic and density functional theory. 

\section{Introduction}
Recent structure-property studies have linked spinodal decomposition to stiffening of bcc delta ferrite, \cite{al2023detecting, al2020detecting,al2020non} this paper sets out a theoretical foundation for the predictability of this phenomenon from quantum elastodynamic first principle, by combining elastodynamic and density functional theory. 


For metals with melting points (MP) above 900°C, which include body-centered cubic (bcc) $\delta$-ferrite (MP 1535°C), face-centered cubic (fcc) γ-austenite (MP 910°C) and stainless steel (ballpark MP range 1400-1530°C, CASS MP 1375-1510°C, 304-type stainless steel (SS304) MP 1400-1450°C), elasticity can be related to plasticity — which is consistent with dislocation theory such that energy of dislocations is stored outside the core in the Hookean part of the lattice. Thermally activated processes may occur during plastic deformation and alter elastic properties. \cite{pugh1954xcii} The cohesive properties of a metal, which can be quantified by lattice elasticity, also related to melting point by their dependence on interatomic potential and electronic structure \cite{chang1966correlation,fine1984elastic, korzhavyi2009electronic, razumovskiy2011first, hosemann2018materials} in ferromagnetic alloys change as a function of concentration. \cite{korzhavyi2009electronic} The origin of the stiffening observed by SAW peak splitting in spinodally decomposed SS304, measured by TGS and simulated in MD for binary FeCr in this paper, is related to an anomalous peculiarity in concentration dependence of elastic constants that leads to a change in the Fermi surface topology. \cite{razumovskiy2011first}

During spinodal decomposition of bcc $\delta$-ferrite, the change in concentration of chromium leads to a change in stiffness in Fe-rich $\alpha$ and Cr-rich $\alpha$′ phases of SS304. One can, therefore, predict plastic embrittlement from elastic stiffening \cite{du2017facile, sdajani2019cassspinodal,al2020non, hofmann2019transient, al2020detecting} in stainless steels undergoing spinodal decomposition, whether thermally-induced and/or radiation-induced. \cite{byun2016thermal,lach2020precipitation} Based on this, one can infer plastic deformation, such as embrittlement, from elastic properties, such as the elastic tensor and moduli. The elastic (Young’s) modulus is directly correlated to surface acoustic wave frequency, which can be measured with picosecond ultrasonic pump-probe spectroscopy. \cite{hofmann2019transient}


%- CASS intro

%- FeCr intro
Fe-rich alloys, dynamically stabilized by other alloying metals, serve as models for Earth's inner core and other systems. To understand the interplay between elastic and magnetic properties, ab initio lattice dynamical calculations have been employed to determine seismic velocities at high temperatures, which are consistent with observations. \cite{dziewonski1981preliminary, oganov2004theoretical,karato1998seismic, wentzcovitch2010first,kadas2011thermo, euroresearchcouncil2014} Additionally, temperature-dependent quantum mechanical calculations of phonon dispersion in bcc Fe phases have been investigated, revealing a dynamically stabilized Fe-rich structure due to phonon-phonon coupling and demonstrating dynamical stability.

Recent studies examining the dependence of elastic properties and free energy on the magnetic moment of Fe atoms provide insights into the interconnectedness of these properties. Understanding the interplay between elastic and magnetic properties is crucial for comprehending the effects of spinodal decomposition and its implications on plastic embrittlement. Spinodal decomposition-induced demixing is observed in various systems, including Earth's inner core, solar cells, and stainless steel used as perovskite solar cell substrates. Based on evidence through secondary scanning electron micrographs, transmission electron micrographs, and visualizations of aged duplex stainless steel welds, these observations highlight the impact of spinodal immiscibility on the seismic, optoelectronic, and mechanical properties of the respective systems.




%- Elastodynamics intro




%- DFT intro





In this work, we employ \emph{ab initio} calculations to explore the fundamental mechanisms at the quantum mechanical level that govern the behavior of materials undergoing spinodal decomposition. By integrating the principles of elastodynamics with those of density functional theory (DFT), we aim to provide a comprehensive understanding of the relationship between microstructural changes and macroscopic physical properties. Our approach enables us to elucidate the connections between atomic-scale phenomena and observable elastic anomalies, offering predictive insights into the mechanical performance of advanced materials in critical applications.